\subsection{Finite genus-one actions and the actual universal-curve transport domain}

For a marked compact oriented surface $S$, let $\Gamma(S)$
be its mapping class group, allowing the specified permutations
of marked points and parametrized boundary components.  Let
$P\Gamma(S)$ be the subgroup fixing all these labels and
fixing each parametrized boundary pointwise.  We use surfaces
of genus zero or one with at most twelve marked points and
boundary components in total.  Stable surfaces are included;
the disc, annulus, and closed torus are also included for
their actual capping maps.  A label stabilizer is always
taken inside the indicated mixed group.

For an integer $m\geq2$, set
\[
 \begin{gathered}
 U_S(m)=\Gamma(S)/\kappa_m(P\Gamma(S)),\qquad
 B(m)=12!(m!)^{(m!)^{90}}.\\
 M_{\rm IV}(b)=\text{the bound (FEC.22), evaluated with
                       $|Q|=\exp(Q)=b$},\\
 L(b)=(M_{\rm IV}(b)!)^{(M_{\rm IV}(b)!)^{45}},\qquad
 N(m)=\max\{B(m),L(B(m))\},\\
 Z_1(m)=\Gamma_{0,5}/\kappa_{N(m)}(\Gamma_{0,5}).
 \end{gathered}
                                                               \tag{FGA.1}
\]
The second line denotes the explicit monotone integer
expression in (FEC.9), (FEC.18), (FEC.20)--(FEC.22),
including the three finite recursions (HCF.11)--(HCF.12)
and (KFU.11).  It does not mean a search for a quotient
with an unspecified property.

\begin{result}{A.4.40}{Proved}{One finite five-point input produces the retained genus-one actions and surface automorphisms}
Every actual gentle pair $F=(\lambda,f)$ in the finite
pentagon locus
\[
      \mathcal P_1(m)=
      \{F\in\widehat{\mathrm{GT}}_0:
                       p(f)=1\text{ in }Z_1(m)\}
                                                               \tag{FGA.2}
\]
gives automorphisms of all the groups $U_S(m)$, compatible
with their actual labelled forgetting, marked capping,
unmarked capping and ordinary sewing maps on the specified
stabilizers.  Quilt and basepath choices give the explicit
inner comparisons below.  The action has the same parameter
$(\lambda,f)$ and the coefficient $c_2(f)$ in the local
formula (TA.8).  The proof of invertibility uses the actual
finite gentle record; it does not assume that the locus
(FGA.2) is a group.

For $1\leq n\leq11$, put $\Pi_n=\pi_1(S_{1,n})$, of
free rank $r=n+1$.  For every specified $s\geq2$ define
\[
 \begin{gathered}
 Q(s,n)=\Pi_n/\kappa_s(\Pi_n),\qquad
 h_s=(s!)^{(s!)^r},\qquad
 w_s=h_s3^{1+h_s(r-1)},\qquad m_{s,n}=w_s^r.
 \end{gathered}
                                                               \tag{FGA.3}
\]
The action supplied by $\mathcal P_1(m_{s,n})$ induces an
actual automorphism $F_{s,n}$ of $Q(s,n)$ such that
\[
 \begin{gathered}
 F_{s,n}(C_i)\sim C_i^\lambda\quad(1\leq i\leq n),\\
 F_{s,n}\,G(s,n)\,F_{s,n}^{-1}=G(s,n),\qquad
 G(s,n)=\operatorname{im}\bigl(
       \Gamma_{1,[n]}\to\operatorname{Out}(Q(s,n))\bigr).
 \end{gathered}
                                                               \tag{FGA.4}
\]
Conjugacy of the peripheral elements in (FGA.4) is
conjugacy inside the actual finite surface group.
In particular, with $n=9$ and
$s=128\cdot3^{1172}$, the multiplier-one action lies
in precisely the finite transport domain (TF.2).

At the two-point torus, its comparison with the actual
$B_4/\langle\omega_4\rangle$ presentation is the pure
gauge (ABR.2)--(ABR.11), evaluated in the retained finite
group.  The full ordinary $J_c$ forgetting factor can
be retained simultaneously as described below.  Equality
of that factor with the different extraction map $d$
on every element of a finite translation normalizer is
not asserted by this result.
\end{result}

\begin{proof}
\par\medskip\noindent\textbf{The explicit finite group bound.}\quad
For a torus with one parametrized boundary, its group is
$B_3$, generated by two twists; its first boundary twist
is already their word $(ab)^6$.  Adding up to eleven
further marked points gives Birman point-pushing kernels
of ranks $2,3,\ldots,12$.  Each additional point can then
be blown up to a parametrized boundary by adjoining at
most its one boundary twist.  Thus the pure groups in
genus one under consideration have generating sets of
size at most
\[
               2+\sum_{j=2}^{12}j+11=90.
                                                               \tag{FGA.5}
\]
This count keeps the first boundary twist inside the
two-generator base group.  Equivalently, capping one
of $b\geq2$ boundary components has kernel generated
by the $b$ simple point-pushing loops and the new
boundary twist; this gives $b(b+3)/2$ generators.
These generators come from the actual evaluation
fibration for a marked point or parametrized disc,
the same maps used in (TA.11)--(TA.13).

In genus zero, the three-boundary group is generated
by its three boundary twists.  Adding the $b$-th
boundary adds at most $b-2$ point-pushing generators
and its boundary twist.  The resulting bound is
$b(b-1)/2\leq66$.  Marked capping gives the punctured
groups as actual quotients of these groups.  The
disc and annulus have respectively zero and one
generator; the closed torus is the actual quotient
$B_3/\langle(ab)^6\rangle$ and has two.
The allowed permutation quotient has order at most
$12!$ in every case.  The diagonal definition of
$\kappa_m$ therefore gives
\[
                  |U_S(m)|\leq B(m).
                                                               \tag{FGA.6}
\]
The kernel is normal in the mixed group because it
is characteristic in its pure normal subgroup.

We will repeatedly use one exact factorization rule.
For a label stabilizer $A$ containing $P\Gamma(S)$,
every finite map $A\to D$ with $|D|\leq m$ kills
$\kappa_m(P\Gamma(S))$: restrict to the pure group,
then compose with the regular representation of $D$
and add fixed letters to degree $m$.
Thus that map is defined on the corresponding
subgroup of $U_S(m)$.  This is a statement about the
actual pure kernel and the specified stabilizer;
it does not require an extension of the homomorphism
to the whole mixed group.

\par\medskip\noindent\textbf{All local pants relations hold in the retained groups.}\quad
Let $Q=U_S(m)$ and evaluate every local quilt weight
in this finite group.  Any homomorphism from the
mixed group $\langle u,w^2\rangle\leq B_3$ into $Q$
has actual image of order at most $B(m)$.
By monotonicity of the explicit bound in A.4.39,
the quotient $Z_1(m)$ maps to each of its required
five-point inputs.  Consequently every element of
(FGA.2) satisfies (TA.2) in every such image.
This includes all four literal mixed maps (TA.3),
on every embedded two-boundary torus.  Conjugating
an embedding conjugates its evaluation in $Q$, so
no infinite list of additional tests is introduced.

Likewise, every local five-A cell is the ordinary
five-point pentagon, possibly with commuting boundary
twists in its arguments.  Thicken the four strands
of $P_4$ to four holes with their constant framings;
this gives the actual homomorphism into its
five-holed-sphere support.  Compose with the support
inclusion and with $\Gamma(S)\to Q$.
Its image has order at most $B(m)$.  The central
cap (FPA.7)--(FPA.9) has order at most that same
bound, so is retained by $\kappa_{N(m)}(\Gamma_{0,5})$.
The zero-character argument there proves the
pentagon in the original finite image, including
the boundary-center coordinate.  Commuting boundary
factors may also be deleted from an argument of
$f$ by the actual direct-product homomorphism used
in (ABR.11), since $f$ has zero exponent sum.

The other local cells require only the genuine gentle
two-cycle and central-product three-cycle identities.
The calculations (TA.3)--(TA.8) are algebraic
equalities under the stated homomorphisms, so remain
valid in $Q$: the six-AS cell uses the four just
established finite mixed identities and this finite
pentagon; the three-S cell cancels for every value
of $c_2(f)$.  A- and S-backtracking, three-A cells,
disjoint supports, and the terminal quilt adjustments
have exactly the checks recorded there.
These are all the cells of the simply connected
pants complex, with the same ordinary presentation
and local supports used in A.6.1
\cite[Theorem~8.2, Definition~7.2 and Claim~8.3]{ET_NS2000}.
No relation at an unretained quotient is being used.

\par\medskip\noindent\textbf{The finite path functor and its universal kernel.}\quad
Write $c=c_2(f)$ and $\nu=(\lambda-1)/2$.
Use the exact quilt weights from A.6.1:
an adjustment $D_a^{N/2}$ has weight $D_a^{N\nu}$,
an A-edge has weight $a^{N\nu}f(b,a)$, and an S-edge
has weight
\[
             a^{N\nu-8c}f(b^2,a^2)b^{8c}(aba)^{2\nu}.
                                                               \tag{FGA.7}
\]
All values in this paragraph are evaluated in $Q$.
The local cell equalities make the product along a
quilt path independent of the path.  Indeed, every
pants loop is a product of the listed ordinary
cells and backtracking, and the free collar
adjustment lattice cancels the terminal adjustment.
The proof of that lattice in A.6.1 cuts each pant
into two hexagons and records the actual collar
displacements; it is unchanged by evaluating the
weights in a finite group.
We therefore obtain finite transports $T_F(Q_1,Q_2)$
with the two literal identities
\[
 \begin{aligned}
 T_F(Q_1,Q_3)&=T_F(Q_1,Q_2)T_F(Q_2,Q_3),\\
 T_F(hQ_1,hQ_2)&=q(h)T_F(Q_1,Q_2)q(h)^{-1}.
 \end{aligned}
                                                               \tag{FGA.8}
\]
Choose a starting quilt $Q_S$ for each surface.
For an ordinary mapping class $\gamma$ set
\[
       \Psi_F(\gamma)=q(\gamma)
                     T_F(Q_S,\gamma^{-1}Q_S)^{-1}.
                                                               \tag{FGA.9}
\]
Substitution of (FGA.8) proves that this is a
homomorphism $\Gamma(S)\to Q$.  Every weight is in
the image of the pure mapping class group, since
it is a word in Dehn twists, with profinite powers
evaluated in that finite group.  Thus (FGA.9)
preserves the original label permutation.

It remains to prove descent through the new kernel;
ordinary geometric functoriality alone would not
prove this.  On the pure subgroup (FGA.9) is a
homomorphism to $P\Gamma(S)/\kappa_m(P\Gamma(S))$.
Compose it with every defining coordinate into
$\mathfrak S_m$.  Each composite kills the source
$\kappa_m$ by its definition.  Those coordinates
are jointly faithful, so (FGA.9) kills that kernel.
It therefore induces an endomorphism of the
specified finite group $U_S(m)$.

\par\medskip\noindent\textbf{Invertibility and composition at the actual finite source record.}\quad
Set
\[
 E=\widehat F_2/\kappa_{B(m)}(\widehat F_2),\qquad e=|E|.
                                                               \tag{FGA.10}
\]
Every evaluation of $f$ in a group of order at most
$B(m)$ factors through $E$.  The order of each
twist occurring in (FGA.7) divides $e$, by applying
the same fact to the cyclic image of a generator.
Thus every finite local weight is determined by
the record
\[
             (\lambda\bmod2e,\ f\bmod\kappa_{B(m)}(F_2),
                                      c_2(f)\bmod e).
                                                               \tag{FGA.11}
\]
This is an actual finite group image of the gentle
group.  The normalized gentle automorphism preserves
the characteristic kernel defining $E$, and its
action there is determined by the first two entries.
The exact third-coordinate law is (TA.1), so
multiplication on the actual records is
\[
  (\lambda,f,c)(\mu,g,c')
       =(\lambda\mu,\ f\beta_F(g),\ c+\lambda c'),
 \qquad
 \beta_F(x)=x^\lambda,\quad\beta_F(y)=f^{-1}y^\lambda f.
                                                               \tag{FGA.12}
\]
The coefficient entry itself is continuous through
the actual index-two Schreier abelian quotient
modulo $e$, of order $2e^3$; its value is not a
freely chosen coordinate.  The image in (FGA.11)
has order at most $\varphi(2e)e^2$.

Composition on an A-edge is the first two-coordinate
law in (FGA.12).  On an S-edge it is exactly
(TA.8), with $c+\lambda c'$ and
$\nu+\lambda\nu'$; adjustments have the same scalar
law.  These are equalities for the finite local
weights along any specified path.  Applying the
already constructed endomorphism to the successive
weights proves the same composition rule along
that whole path.  In particular the $k$-th power
of (FGA.9) has on every chosen path the record of
the actual gentle power $F^k$.  This assertion does
not require a separate pentagon premise for $F^k$.

Some positive power of the actual record (FGA.11)
is the identity.  For that record $\nu=0\pmod e$,
$c=0\pmod e$, and every value of $f$ is one, so
each local path weight is one.  Formula (FGA.9)
then fixes every ordinary mapping class at the
finite quotient.  Hence that power of the
endomorphism is the identity, proving invertibility.
For a uniform integer one may use
$\bigl(\varphi(2e)e^2\bigr)!$.
This argument uses a finite group of actual source
records, not a claimed subgroup property of (FGA.2).

\par\medskip\noindent\textbf{The actual sewing and forgetting maps.}\quad
An embedded quilt path remains the same local
path after sewing.  The two boundary twists sewn
to one seam have positive images whose exponents
add.  Thus every weight maps to the sewn weight.
Under a cap, the adjacent pant becomes an annulus;
the two surviving collar displacements add, as
in the literal strip calculation (TA.11).
An A-move meeting the cap has commuting replacement
twists after capping, so its commutator coefficient
is one and its remaining weight is precisely
that collar adjustment.  An S-move is unchanged
unless its entire one-holed-torus support is capped.
That last case is the explicit quotient
$B_3\to\operatorname{SL}_2(\mathbb Z)$, whose
kernel is the boundary twist $(ab)^6$.

The kernel for marked capping is generated by the
capped boundary twist.  For unmarked capping it
is generated by that twist and the point-pushing
bounding-pair products.  A common quilt containing
both curves gives, in the finite group,
\[
 \begin{aligned}
 \Psi_F(D_cD_d^{-1})
     &=t(D_cD_d^{-1})^\lambda t^{-1},\\
 \Psi_F(D_\delta)&=t_\delta D_\delta^\lambda t_\delta^{-1}.
 \end{aligned}
                                                               \tag{FGA.13}
\]
These are the actual generators from the disc
evaluation fibration.  Their images are killed
after the specified cap, so the finite homomorphism
descends through the ordinary capping quotient.
On its own universal pure kernel it descends again
by the same $\mathfrak S_m$ coordinate argument.
The finite inverse power proves preservation with
equality of the relevant finite kernel images.
Successive caps add the same collar displacements,
so their order gives the same map.  This proves
forgetting on the specified label stabilizers.

Every ordinary pure sewing, relabelling or forgetting
map carries source $\kappa_m$ into target $\kappa_m$
by composition with its target coordinates in
$\mathfrak S_m$.  Thus the preceding comparisons
are comparisons of the actual finite quotient maps.
Choosing quilts by the actual capping rule makes
each stated forgetting square literal.  With other
starting quilts, its inner element is the explicitly
defined path weight $T_F$ between those two quilts.
No such element is silently set to one.

\par\medskip\noindent\textbf{A finite center-killing cover of the punctured torus group.}\quad
We now construct the surface action rather than
infer it from an action on an inner automorphism
quotient.  Let $1\leq n\leq11$, $r=n+1\geq2$,
and $Q_0=Q(s,n)$.  Put
\[
 K=\kappa_s(\Pi_n),\qquad
 Q_1=\Pi_n/(K'K^3),\qquad h=|Q_0|.
                                                               \tag{FGA.14}
\]
All kernels are characteristic in the actual free
surface group, and
\[
 |Q_1|=h3^{1+h(r-1)}\leq w_s,
 \qquad Z(Q_1)\longrightarrow Q_0\text{ has image }1.
                                                               \tag{FGA.15}
\]
Here is the exact second assertion.  The regular
graph cover corresponding to $K$ has first integral
homology of rank $1+h(r-1)$.  A nonidentity deck
element has trace one on this homology: its traces
on vertices and on oriented edges are both zero,
and its trace on $H_0$ is one.  The identity has
the strictly larger trace $1+h(r-1)$.
The deck action is therefore faithful over the
integers.  A finite-order integral matrix congruent
to the identity modulo three is the identity;
the elementary binomial-valuation proof is the
one retained in (DR.6).  Hence the same deck action
is faithful on $K/K'K^3$.  A central element of
$Q_1$ acts trivially there, so its deck image in
$Q_0$ is trivial.  This proves (FGA.15).
In particular there is a genuine surjection
\[
       \operatorname{Inn}(Q_1)=Q_1/Z(Q_1)
                             \longrightarrow Q_0.
                                                               \tag{FGA.16}
\]
The free, punctured scope is essential to this
argument.  It is not an assertion about the
fundamental group of an unpunctured closed torus.

\par\medskip\noindent\textbf{The actual Birman subgroup and its characteristic output.}\quad
Let $\Gamma_+$ be the subgroup of
$\Gamma_{1,[n+1]}$ fixing the last, newly added
point.  Its ordinary Birman sequence is
\[
 1\longrightarrow\Pi_n\longrightarrow\Gamma_+
     \longrightarrow\Gamma_{1,[n]}\longrightarrow1.
                                                               \tag{FGA.17}
\]
Conjugation on this named kernel is its based
surface action.  It gives the actual finite map
$\rho_1:\Gamma_+\to\operatorname{Aut}(Q_1)$,
whose image has order at most $|Q_1|^r\leq m_{s,n}$.
The label-stabilizer factorization after (FGA.6)
shows that it factors through the subgroup $U_+$
of $U_{1,n+1}(m_{s,n})$ fixing that last label.
Let $P$ be the actual image of $\Pi_n$ in $U_+$.
Its image under $\rho_1$ is exactly
$\operatorname{Inn}(Q_1)$.  Thus (FGA.16) supplies
a surjection $P\to Q_0$, whose composite with
$\Pi_n\to P$ is the originally specified quotient
to $Q_0$.  In particular
\[
 \ker(\Pi_n\to P)\subseteq\kappa_s(\Pi_n),
 \qquad P/\kappa_s(P)\xrightarrow{\sim}Q_0.
                                                               \tag{FGA.18}
\]
For the second equality, every homomorphism
$\Pi_n\to\mathfrak S_s$ factors through $P$ by
the first inclusion, and every homomorphism from
$P$ pulls back to such a homomorphism.  Their
kernel intersections therefore correspond exactly.
This proves the displayed isomorphism by its
actual quotient maps.

The automorphism of $U_{1,n+1}(m_{s,n})$ preserves
$U_+$ because it preserves label permutations.
It also preserves $P$.  Indeed the Birman kernel
is generated by the point pushes along a simple
free surface basis.  Each is the actual product
$D_cD_d^{-1}$ around the pant traced by that loop,
after killing the new boundary twist by marked
capping.  Equation (FGA.13) sends its image to a
conjugate of its $\lambda$th power, inside the
normal subgroup $P$.  This gives containment;
invertibility on the finite group gives equality.
Consequently the automorphism on $P$ descends to
the characteristic quotient (FGA.18), giving the
claimed actual $F_{s,n}\in\operatorname{Aut}(Q_0)$.
It is not merely an automorphism of
$\operatorname{Inn}(Q_0)$.

For each named peripheral loop $C_i$, its point
push is the twist about the curve enclosing the
new point and the $i$th old point, with the
consistent orientation from (FGA.17).
The twist formula gives a conjugate of its
$\lambda$th power.  The transport conjugator is
in the finite image of the pure mapping class
group, as noted after (FGA.9).  Choose its ordinary
pure representative.  Its based surface action
sends $C_i$ to a conjugate of $C_i$ inside
$\Pi_n$.  This proves the conjugacy assertion
of (FGA.4) in $Q_0$ itself.

Finally, $K/(K'K^3)$ is exactly $\kappa_s(Q_1)$:
the same kernel-intersection argument applies to
$\Pi_n\to Q_1\to Q_0$.  Thus $Q_0$ is also a
characteristic quotient of $Q_1$, and the retained
$\rho_1$ projects to the actual based geometric
action on $Q_0$.  For $g\in U_+$, covariance of
conjugation on $P$, followed by (FGA.18), gives
\[
 F_{s,n}\rho_0(g)F_{s,n}^{-1}
                       =\rho_0(\Psi_F(g)).
                                                               \tag{FGA.19}
\]
The right side belongs to the same actual geometric
image.  The inverse automorphism gives equality
of those finite images; passing to outer classes
proves the second assertion of (FGA.4).
Compatible base quilts make its action on the
base of (FGA.17) the already constructed
finite action on $\Gamma_{1,[n]}$.

\par\medskip\noindent\textbf{The two-point braid gauge is retained at this level.}\quad
On the marked two-boundary torus, use the actual
chain $a,b,d$ and the two adapted moves of A.4.26.
The finite actions just constructed have the three
images (ABR.2), because those are precisely their
local A- and S-edge weights.  All four mixed
coefficient evaluations needed in (TA.4), including
the $(d,V)$ evaluation (ABR.4), are already forced
by the finite input (FGA.2).
The proof of (ABR.7)--(ABR.11) is therefore an
identity in the retained finite group, with the
single specified pure gauge
\[
                   g=a^{8c_2(f)}V^{-4c_2(f)},
                   \qquad V=(aba)^2.
                                                               \tag{FGA.20}
\]
It identifies the three images with the standard
four-strand braid action.  Capping gives the
actual quotient $B_4/\langle\omega_4\rangle$
and its index-twenty-four pure subgroup
$P_4/\langle\omega_4\rangle=\Gamma_{0,5}$,
with the same pure inner comparison.  Thus this
finite branch-order comparison follows from the
specified finite pentagon input and its exact
mixed evaluations, without assuming a full GT
parameter at the intermediate step.

\par\medskip\noindent\textbf{Retaining the full ordinary compact factor.}\quad
There is a simultaneous input for the nine-point
transport and the full factor $J_c$, not only its
signed period.  Put
\[
 s_9=128\cdot3^{1172},\qquad s_c=|\mathcal P_c|=32\cdot3^{14},
 \qquad
 m_{\rm CM}=\max\{m_{s_9,9},m_{s_c,1},|J_c|\},
                                                               \tag{FGA.21}
\]
where the two $m$ values mean (FGA.3) with the
displayed pair $(s,n)$.
The larger universal level maps to both smaller
ones, and the finite path formulas intertwine
these projections.  For $F\in\mathcal P_1(m_{\rm CM})$
we thus have the two actual surface automorphisms
$F_9$ and $F_\Sigma$ on their characteristic inputs.

The latter descends to $\mathcal P_c$.
Here is the precise peripheral invariance at this
finite input.  The quaternion kernel $N_8$ is
recovered as the unique kernel of an epimorphism
to $Q_8$: its twenty-four generating pairs are
one automorphism orbit.  The subgroup $K_3$ is
the kernel of ordinary abelian reduction modulo
three, and $T=N_8\cap K_3$.  Their images, the
verbal subgroup $L_8=\gamma_3N_8\,N_8^3$, and
$T'$ are therefore preserved.  The subgroup
$D_\partial$ is the normal closure of the square
of the named peripheral commutator, and the
peripheral power rule preserves that closure
with equality because $\lambda$ is a unit.
Finally the abelian modulo-twelve kernel is
characteristic.  These operations recover exactly
\[
 R_c=(L_8T'D_\partial)\cap\Sigma'\Sigma^{12}
                                                               \tag{FGA.22}
\]
inside the retained characteristic input: its
kernel is already contained in $R_c$ by the
regular representation of $\mathcal P_c$.
Intersections commute with this projection
because its kernel is contained in both factors.
Thus $F_\Sigma$ gives an actual automorphism
$F_c$ of $\mathcal P_c$.
Covariance (FGA.19) on the one-point torus now
normalizes its actual ordinary image $J_c$.

Let $M=\Gamma_{1,[9]}$ and let $M_0$ fix the
distinguished zero label.  Retain the actual map
\[
 j_0:M_0\xrightarrow{\text{forget eight labels}}
            \Gamma_{1,1}\longrightarrow J_c.
                                                               \tag{FGA.23}
\]
It factors through the specified subgroup of the
universal finite level, by (FGA.6).
Choose the base quilts by successive capping;
the two instances of (FGA.19) and forgetting
compatibility give the exact square
\[
                    j_0\Psi_F=\operatorname{Ad}(F_c)j_0.
                                                               \tag{FGA.24}
\]
With a different starting quilt on the one-point
target the right side is conjugated by its named
element of $J_c$, the image of the path weight
in (FGA.8).  This leaves the same kernel and is
the recorded inner ambiguity.
Consequently the entire induced nine-coset map
$\operatorname{Ind}_{M_0}^{M}(j_0)$ is retained
and its kernel is preserved: it is the intersection
of the nine actual conjugates of $\ker j_0$,
and $\Psi_F$ preserves the label stabilizer and
the kernel in (FGA.24).
Joining that map to the surface action gives the
actual finite geometric image
\[
 G^\dagger=\operatorname{im}\left(
  M\longrightarrow
  G\times\operatorname{Ind}_{M_0}^{M}(J_c)\right),
 \qquad G\leq\operatorname{Out}(Q_{\rm out}),
                                                               \tag{FGA.25}
\]
on the multiplier-one domain of A.4.25.
Both coordinates and their source are kept in
this joint image; no independence or direct-product
surjectivity is asserted.

\par\medskip\noindent\textbf{The exact remaining comparison and multiplier transfer.}\quad
On the zero-label subgroup of the finite normalizer
$N_{G^\dagger}(H)$ there are two named homomorphisms
to $J_c$: the identity-coset component of (FGA.23),
and $d$ from (TF.4) after projection to $G$.
They agree on the actual ordinary deck-normalizer
source by its geometric covering/forgetting map.
The present construction does not extend that
equality to every element of the finite normalizer.
The explicit discrepancy of A.4.42 shows why
this further recognition step cannot be supplied
by merely adjoining the nine solvable forgetting
coordinates to the existing homology output.

The restriction to multiplier one in the transport
fusion theorem does not discard the other actual
gentle multipliers.  For an actual $F$ in (FGA.2),
choose an actual arithmetic element $\gamma$ with
$\lambda(\gamma)=\lambda(F)$, using the surjective
cyclotomic character.  Its full GT action gives
an invertible finite functor at the same level.
Compose the already constructed functor for $F$
with its inverse.  The local composition calculation
(FGA.12) makes the resulting parameter exactly the
actual gentle product $F\gamma^{-1}$, of multiplier
one, on every retained path and quotient.
This constructs its coherent finite action directly;
it does not assert that $F\gamma^{-1}$ belongs to
the original finite pentagon locus.
The supplied character is zero on arithmetic
elements and is a homomorphism, so
$\Delta(F\gamma^{-1})=\Delta(F)$.
Thus any completed finite comparison on this
coherent multiplier-one domain transfers to all
actual multipliers through these explicit
composition maps.  The unresolved comparison just
identified, and the actual-source pentagon-shadow
census, are not replaced by this transfer.
\end{proof}
